% !TeX root = ../main.tex
% ---------------------------------------------------------------------------
% Advanced Numerical Methods
% Scope extension: extends Chapter 8 ( Chapter~\ref{ch:numerical-modelling}, 2D continuum FEM in RS2).
% Key sources: Potts & Zdravkovic; Jing (2003); Cundall & Hart; Barla on
% time-dependent modelling; ITASCA and Rocscience documentation.
% ---------------------------------------------------------------------------
\chapter{Advanced Numerical Methods}
\label{ch:advanced-numerical}


\section{Scope and limits of numerical analysis}
\label{sec:scope-and-limits-of-numerical-analysis}

\subsection{The limits of 2D continuum analysis}
\label{sec:limits-2d}
\scopenote{Plane strain assumptions, the stress relaxation factor as a proxy
for three-dimensional face effects, and the problem classes where that proxy
fails.}

\subsection{Constitutive models beyond Mohr--Coulomb}
\label{sec:constitutive-models}
\scopenote{Hardening soil and small-strain stiffness models for urban
settlement; strain-softening and cohesion-weakening friction-strengthening
(CWFS) for brittle rock; creep and viscoplastic models for squeezing ground.}

\subsection{Selecting and parameterising a constitutive model}
\label{sec:model-selection}
\scopenote{Matching model complexity to available data. The common error of
adopting an advanced model with parameters that were never measured.}

\section{The family of numerical methods}
\label{sec:the-family-of-numerical-methods}

\subsection{The family of numerical methods}
\label{sec:numerical-family}
\scopenote{Boundary element, finite element, finite difference and distinct
element methods compared by what each discretises, the problems each suits,
and their differing treatment of far-field boundaries. After Brady \& Brown
Ch.~6.}

\subsection{Boundary element methods}
\label{sec:bem}
\scopenote{Discretising only the excavation surface: the efficiency advantage
for elastic problems in an infinite medium, the natural handling of far-field
conditions, and the restriction to linear homogeneous behaviour. Still the
quickest way to map stress around a complex excavation shape.}

\subsection{Linked and hybrid schemes}
\label{sec:linked-methods}
\scopenote{Coupling a boundary element far field to a finite or distinct
element near field, so that nonlinearity is modelled only where it occurs.}

\section{Discontinuum and hybrid methods}
\label{sec:discontinuum-and-hybrid-methods}

\subsection{Discontinuum modelling}
\label{sec:discontinuum}
\scopenote{Distinct element method, block generation, joint constitutive
behaviour, and the problems where explicit joints change the answer rather
than merely refining it. UDEC and 3DEC in practice.}

\subsection{Block theory and key block analysis}
\label{sec:block-theory}
\scopenote{Goodman \& Shi's removability theorem as a rapid deterministic
screen for structurally controlled failure, and its relationship to the wedge
analyses of Chapter~\ref{ch:rock-reinforcement}.}

\subsection{Hybrid and fracture-propagation methods}
\label{sec:fdem}
\scopenote{Combined finite--discrete element (FDEM) modelling of fracture
initiation and propagation; applications to spalling, blast damage and
fragmentation.}

\section{Three-dimensional and staged analysis}
\label{sec:three-dimensional-and-staged-analysis}

\subsection{Three-dimensional modelling of the advancing face}
\label{sec:3d-modelling}
\scopenote{Where 3D is genuinely required: junctions, cross-passages, portals,
cavern sequences, twin tunnel interaction and fault crossings. Cost and mesh
management.}

\subsection{Simulating construction sequence}
\label{sec:construction-sequence}
\scopenote{Staged excavation, timing of support activation, shotcrete ageing,
segment ring installation and grouting, and the sensitivity of results to
sequence assumptions.}

\section{Coupled and multi-physics analysis}
\label{sec:coupled-and-multi-physics-analysis}

\subsection{Coupled hydro-mechanical analysis}
\label{sec:coupled-hm}
\scopenote{Consolidation and steady-state seepage coupling, drainage boundary
conditions on the lining, and the time-dependent settlement this produces.}

\subsection{Thermal, dynamic and multi-physics problems}
\label{sec:multiphysics}
\scopenote{Thermo-hydro-mechanical analysis for repositories and energy
geostructures; dynamic analysis for blasting and seismic loading
(Chapter~\ref{ch:seismic-design}).}

\section{Calibration and quality assurance}
\label{sec:calibration-and-quality-assurance}

\subsection{Calibration, back-analysis and model updating}
\label{sec:back-analysis}
\scopenote{Using monitoring data to update the model during construction, the
non-uniqueness of back-analysis, and its formal role in the observational
method.}

\subsection{Verification, validation and quality assurance}
\label{sec:model-qa}
\scopenote{Mesh convergence, benchmark problems, sanity checks against
closed-form solutions, independent review, and the professional obligation not
to present unverified output as design.}
